% BHU PYQ -- Manually transcribed & error-corrected
% Paper: BCH-223 : Specialised Accounts II
% Exam: B.Com. (Hons.) Semester IV Examination, 2022-23

\documentclass[11pt,a4paper]{article}
\usepackage[a4paper,top=2.5cm,bottom=2.5cm,left=2.8cm,right=2.8cm,headheight=14pt]{geometry}
\usepackage{amsmath,amssymb}
\usepackage[T1]{fontenc}
\usepackage[utf8]{inputenc}
\usepackage{lmodern,microtype,fancyhdr,enumitem,hyperref,lastpage,parskip}

\hypersetup{
    colorlinks=true,
    urlcolor=black,
    linkcolor=black,
    pdftitle={BCH-223: Specialised Accounts II -- BHU B.Com. (Hons.) Sem IV 2022-23}
}

\pagestyle{fancy}
\fancyhf{}
\renewcommand{\headrulewidth}{0.4pt}
\renewcommand{\footrulewidth}{0.4pt}
\fancyhead[L]{\small   \textit{Banaras Hindu University}}
\fancyhead[R]{\small   \textit{BCH-223: Specialised Accounts II}}
\fancyfoot[C]{\small Page  \thepage\ of \pageref{LastPage}}


\newlist{parts}{enumerate}{2}
\setlist[parts,1]{label=(\alph*),leftmargin=*,itemsep=5pt,topsep=3pt}
\setlist[parts,2]{label=(\roman*),leftmargin=*,itemsep=3pt,topsep=2pt}

\newcommand{\pts}[1]{\hfill   \textit{[#1]}}


\begin{document}


\begin{center}
  {\large\bfseries BANARAS HINDU UNIVERSITY}\\[2pt]
  {
\normalsize B.Com. (Hons.) --- Semester IV Examination, 2022-23}\\[6pt]
  \rule{0.8\linewidth}{0.5pt}\\[6pt]
  {\large\bfseries Paper: BCH-223: Specialised Accounts II}\\[4pt]
  {
\normalsize Time: 3 Hours\hfill Maximum Marks: 70}
\end{center}

\rule{\linewidth}{0.4pt}

\smallskip



\noindent   \textit{Instructions: Attempt all questions. All questions carry equal marks (14 marks each).}

\bigskip




\noindent   \textbf{Question 1.} \\
Explain the various schedules which are prepared by a Commercial Bank for Balance Sheet and give their specimen.\pts{14}

\centerline{   \textbf{OR}}

From the following information, prepare Profit \& Loss Account of Ram Bank Ltd. for the year ended 31st March, 2023. Also give necessary schedules:\pts{14}


\begin{center}
\renewcommand{\arraystretch}{1.2}
\small

\begin{tabular}{|l|r|r|}
\hline
   \textbf{Particulars} &   \textbf{2021-22 (Rs. 000)} &   \textbf{2022-23 (Rs. 000)} \\
\hline
Interest and Discount & 2,045 & 1,427 \\
Income from Investments & 112 & 114 \\
Interest on Balances with RBI & 177 & 155 \\
Commission & 712 & 722 \\
Profit on sale of Investments & 122 & 12 \\
Interest on Deposits & 822 & 612 \\
Interest paid to RBI & 147 & 127 \\
Payments to and provisions for Employees & 855 & 727 \\
Rent & 179 & 158 \\
Printing and Stationery & 212 & 147 \\
Advertisement & 98 & 112 \\
Depreciation & 98 & 98 \\
Directors' fees & 212 & 148 \\
Auditors' fees & 110 & 110 \\
Law charges & 152 & 50 \\
Postage, Telephones & 62 & 48 \\
Repairs and Maintenance & 52 & 42 \\
Insurance & 66 & 57 \\
\hline
\end{tabular}
\end{center}

   \textbf{Other Information:} The following items are already adjusted with 'Interest and Discount' (Cr.):

\begin{itemize}[leftmargin=*]
  \item Rebate on bills discounted (31.03.2023): Rs. 55,000
  \item Loss on sale of Investments: Rs. 12,000
  \item Provision for Taxation: Rs. 148,000
  \item Provision for Doubtful Debts: Rs. 92,000
\end{itemize}
   \textbf{Appropriations:}

\begin{parts}
  \item 25\% of profit is transferred to Statutory Reserve.
  \item 5\% of profit is transferred to Revenue Reserve.
\end{parts}

\medskip\rule{\linewidth}{0.2pt}

\pagebreak




\noindent   \textbf{Question 2.} \\
Following are the balances in the books of a Railway Company for the year ended 31st March, 2023:\pts{14}


\begin{center}
\renewcommand{\arraystretch}{1.2}
\small

\begin{tabular}{|l|r||l|r|}
\hline
   \textbf{Balances (Dr.)} &   \textbf{Amount (Rs.)} &   \textbf{Balances (Cr.)} &   \textbf{Amount (Rs.)} \\
\hline
Railway Lines & 24,25,000 & Share Transfer Fees & 75,000 \\
Workshop Machinery & 6,00,000 & Discount on Purchase & 50,000 \\
Dead Stock & 2,00,000 & Credit Balance of Net Revenue A/c & 4,50,000 \\
Motor Boats & 1,75,000 & General Interest Receipts & 50,000 \\
Land, Stations \& Building & 4,50,000 & Rent on Leased Lines & 65,000 \\
Engines, Coaches \& Wagons & 8,00,000 & Debenture Interest Outstanding & 1,50,000 \\
Maintenance of Lines \& Stations & 2,00,000 & Cash in hand and with Bank & 5,25,000 \\
Taxes and Rates & 62,500 & Investments & 12,97,500 \\
Locomotive Power & 3,50,000 & Traffic A/c due to the Company & 3,00,000 \\
Compensation & 1,00,000 & Spares Stock & 75,000 \\
General Charges & 50,000 & Reserve Fund & 6,00,000 \\
Traffic Expenses & 2,00,000 & Fire Insurance Fund & 1,55,000 \\
Mileage and Demurrage & 30,000 & Debts due to Other Companies & 75,000 \\
Carriage and Wagons Repair & 75,000 & Sundry Creditors & 1,75,000 \\
Shipping Expenses & 20,000 & Outstanding Dividends and Interest & 50,000 \\
Passenger's Fare & 5,75,000 & Equity Share Capital & 26,25,000 \\
Parcels Freight & 1,75,500 & Preference Share Capital & 14,25,000 \\
Mails Freight & 19,500 & Debentures & 8,50,000 \\
Merchandise Freight & 4,00,000 & Premium on the Issue of Shares & 1,50,000 \\
Minerals Freight & 2,50,000 & & \\
\hline
\end{tabular}
\end{center}

In the year ending on 31st March, 2023, the following additions to fixed assets were made by the company:

\begin{itemize}[leftmargin=*]
  \item Railway Lines: Rs. 3,00,000
  \item Workshop Machinery: Rs. 1,00,000
  \item Engines, Wagons and Coaches: Rs. 2,00,000
  \item Land, Building and Stations: Rs. 75,000
\end{itemize}
On the basis of the above information, prepare the Revenue Account, Net Revenue Account, Capital Account, and General Balance Sheet of the Company by considering the balancing figure.

\centerline{   \textbf{OR}}

Point out the concept and features of Double Account System. How is it different from General Accounting System?\pts{14}

\medskip\rule{\linewidth}{0.2pt}

\pagebreak




\noindent   \textbf{Question 3.} \\
What do you understand by Statement of Affairs? How is it prepared? Give a proforma of Deficiency Account.\pts{14}

\centerline{   \textbf{OR}}

Radhey Shyam commenced business on 1st January, 2019 with a capital of Rs. 67,200. His profits for five years were Rs. 21,600 and his drawings averaged Rs. 4,000 per annum. On 31st December, 2023 an order of insolvency was passed with the following details and having a balancing figure. Prepare the Statement of Affairs and Deficiency Account:\pts{14}

\begin{parts}
  \item Unsecured Creditors: Rs. 40,000
  \item Mortgage on Plant: Rs. 8,000
  \item Creditors partly secured (security being life policy of Rs. 8,000 on which premium was paid personally): Rs. 24,000
  \item Salaries of Naveen for 3 months: Rs. 420
  \item Wages of Ram for 6 months: Rs. 180
  \item Landlord's rent for 2 months: Rs. 160
  \item Owing to Government: Rs. 200
  \item Bills Receivable discounted and expected to rank: Rs. 6,400
  \item Plant, estimated to realize Rs. 40,000: Rs. 80,000
  \item Machinery, estimated to realize Rs. 4,000: Rs. 16,000
  \item Book debts:
    
\begin{parts}
      \item Good: Rs. 12,000
      \item Doubtful (Estimated to realize Rs. 1,200): Rs. 2,000
      \item Bad: Rs. 10,000
    \end{parts}
  \item Furniture, estimated to realize Rs. 700: Rs. 1,600
  \item Stock, estimated to realize Rs. 11,100: Rs. 16,000
  \item Cash in Hand: Rs. 2,160
\end{parts}

\medskip\rule{\linewidth}{0.2pt}




\noindent   \textbf{Question 4.} \\
What is meant by Insurance? Prepare with imaginary figures the Revenue Account and Balance Sheet of a General Insurance Company.\pts{14}

\centerline{   \textbf{OR}}

What are the main types and principles of Insurance? Prepare a fire Revenue Account. How are 'net claims' ascertained?\pts{14}

\medskip\rule{\linewidth}{0.2pt}




\noindent   \textbf{Question 5.} \\
Write an explanatory note on Government Accounting.\pts{14}

\centerline{   \textbf{OR}}

Write lucid notes on any two of the following:\pts{7+7}

\begin{parts}
  \item C \& AG (Comptroller and Auditor General)
  \item Slip System
  \item Preferential creditors in Insolvency Accounts
  \item Re-insurance
\end{parts}

\vfill
\rule{\linewidth}{0.4pt}

\begin{center}\small
  \href{https://bhu-exam.vercel.app/}{Click here to visit our website}\\[4pt]
  \href{https://me-aryan.vercel.app/}{Click here to visit our contributor}\\[4pt]
  \href{https://quashan-me.vercel.app/}{Click here to visit our contributor}
\end{center}

\end{document}
